\chapter{Evaluation and Results}

\begin{table}[htbp]
    \centering
    \caption{Supervised Machine Learning Model Performance Metrics without Optuna Study}
    \label{tab:model_performance_supervised}
    \scriptsize
    \setlength{\tabcolsep}{3pt}

    \begin{tabular}{llccccccc}
        \toprule
        \textbf{Model} & \textbf{Test Dataset} &
        \textbf{Accuracy} & \textbf{Precision} & \textbf{Recall} &
        \textbf{F1} & \textbf{Macro F1} & \textbf{ROC-AUC} & \textbf{MCC} \\
        \midrule

        \multirow[c]{2}{*}{Random Forest}
        & CSE-CIC-IDS2018 & 0.9188 & 0.9239 & 0.2043 & 0.3346 & 0.6457 & 0.8384 & 0.4130 \\
        & UNSW-NB15       & 0.9743 & 0.2830 & 0.0002 & 0.0003 & 0.4936 & 0.7613 & 0.0063 \\

        \midrule

        \multirow[c]{2}{*}{XGBoost}
        & CSE-CIC-IDS2018 & 0.9309 & 0.9799 & 0.3150 & 0.4768 & 0.7199 & 0.9153 & 0.5346 \\
        & UNSW-NB15       & 0.9752 & 0.6533 & 0.0739 & 0.1328 & 0.5601 & 0.7829 & 0.2142 \\

        \midrule

        \multirow[c]{2}{*}{Deep Neural Network}
        & CSE-CIC-IDS2018 & 0.8932 & 0.1393 & 0.0133 & 0.0244 & 0.4839 & 0.5346 & 0.0129 \\
        & UNSW-NB15       & 0.8820 & 0.0492 & 0.1958 & 0.0787 & 0.5078 & 0.8107 & 0.0501 \\

        \bottomrule
    \end{tabular}
\end{table}
All the supervised models generated accuracy greater than 88\%. Random Forest and XGBoost yielded Precision score above 90\% for the 
CSE-CIC-IDS2018 dataset. XGBoost turned out to be the best performing model while evaluated on the same dataset with Recall Score of 0.31, F1-Score of 0.47, 
MCC of 0.53 compared to Recall Score of 0.20, F1-Score of 0.33 and MCC of 0.41 generated by Random Forest. These scores achieved are better than the scores 
achieved from previous studies \ref{tab:evaluation-results}. The performance of these models dropped significantly when evaluated on UNSW-NB15 dataset. 
The only signicant metric being MCC which is 0.21 for XGBoost. The primary reason for the performance drop could be due to variation in the types of 'Attack' 
in the dataset \ref{tab:class_distribution}. The 'Attack' labels that are present in the train dataset CIC-IDS-2017 are not present in the test dataset UNSW-NB15 
and vice-versa. The study of \cite{Chandra2025} have highlighted the models' robustness in cross-dataset generalization for network intrusion detection while testing 
on this particular dataset. Also, the methodology used are in contrast to what we have adopted. 

The biggest disappointment above all is the result from DNN architecture. The model yielded accuracy above 88\% for both datasets, however the results for our 
preferred metrics are negligible/avoidable. As, DNN relies on the depth of the layers to learn complex features and their patterns, huge class-imbalance and the depth
of layers used could have been the reason for the poor performance. We have only used 4 hidden layers to detect attacks from the network traffic. Further more, we 
conducted Optuna study twice to tune the hyperparameters and used them to check whether the results improve, there was no significant change in the results of 
Random Forest and XGBoost, however we saw improvement in the performance of DNN. But tuning the hyperparameters showed inconsistent influence over the test datasets 
CSE-CIC-IDS2018 and UNSW-NB15.  

The following are the hyperparameters we used for our Optuna Studies.

\begin{table}[!ht]
    \centering
    \caption{Hyperparameters for DNN Optuna Studies}  
    \label{tab:Hyperparameters-details DNN}
    {\fontsize{9}{11}\selectfont
    \renewcommand{\arraystretch}{1.25}
    \begin{tabularx}{\textwidth}{
        >{\raggedright\arraybackslash}p{3.5cm}
        >{\centering\arraybackslash}p{2.0cm}
        >{\raggedright\arraybackslash}X
    }

        \toprule

        \textbf{Hyperparameter} &
        \textbf{Search Space} &
        \textbf{Purpose} \\

        \midrule

        n\_layers &
        2, 6 &
        Controls network depth, learns complex features and balances overfitting\\

        units per layer&
        [512, 256, 128, 64] &
        Controls width of hidden layer and checks the respresentational capacity at each layer \\

        learning\_rate &
        [1e-4, 1e-2] &
        Controls the step size during gradient descent \\

        dropout\_rate &
        0.1, 0.5 &
        Deactivate the neurons during training to prevent overfitting and overall regularizes the network \\
        
        batch\_size &
        [512, 1024, 2048, 4096] &
        Prevents memory crashes, training speed and gradient noise by controlling the number of samples processed \\

        \bottomrule
    \end{tabularx}
    }
\end{table}
\FloatBarrier

\begin{table}[htbp]
    \centering
    \caption{DNN Performance Metrics with Optuna Study -1}
    \label{tab:DNN_perofrmance_metrics_optuna1}
    \scriptsize
    \setlength{\tabcolsep}{3pt}

    \begin{tabular}{llccccccc}
        \toprule
        \textbf{Model} & \textbf{Test Dataset} &
        \textbf{Accuracy} & \textbf{Precision} & \textbf{Recall} &
        \textbf{F1} & \textbf{Macro F1} & \textbf{ROC-AUC} & \textbf{MCC} \\
        \midrule

        \multirow[c]{2}{*}{Deep Neural Network}
        & CSE-CIC-IDS2018 & 0.9001 & 0.4135 & 0.0005 & 0.0010 & 0.4742 & 0.5954 & 0.0114 \\
        & UNSW-NB15       & 0.9759 & 0.8357 & 0.0794 & 0.1449 & 0.5664 & 0.8096 & 0.2532 \\
        \bottomrule
    \end{tabular}
\end{table}
\FloatBarrier

The best hyperprameters after Optuna Study 1 were:
\[
\{n\_{\mathrm{layers}},\, \mathrm{units},\, \mathrm{dropout\_rate},\, 
\mathrm{learning\ rate},\, \mathrm{batch\_size}\}
=
\{4,\,512,\,0.179447,\,0.001116,\,512\}.
\]

\begin{table}[htbp]
    \centering
    \caption{DNN Performance Metrics with Optuna Study -2}
    \label{tab:DNN_performance_metrics_optuna2}
    \scriptsize
    \setlength{\tabcolsep}{3pt}

    \begin{tabular}{llccccccc}
        \toprule
        \textbf{Model} & \textbf{Test Dataset} &
        \textbf{Accuracy} & \textbf{Precision} & \textbf{Recall} &
        \textbf{F1} & \textbf{Macro F1} & \textbf{ROC-AUC} & \textbf{MCC} \\
        \midrule

        \multirow[c]{2}{*}{Deep Neural Network}
        & CSE-CIC-IDS2018 & 0.9058 & 0.9824 & 0.0584 & 0.1102 & 0.5302 & 0.5470 & 0.2274 \\
        & UNSW-NB15       & 0.9360 & 0.0025 & 0.0038 & 0.0030 & 0.4850 & 0.7923 & -0.0293 \\
        \bottomrule
    \end{tabular}
\end{table}

The best hyperprameters after Optuna Study 2 were: 
\[
\{n\_{\mathrm{layers}},\, \mathrm{units},\, \mathrm{dropout\_rate},\,
\mathrm{learning\_rate},\, \mathrm{batch\_size}\}
=
\{4,\,512,\,0.111666,\,0.000527,\,512\}.
\]

\begin{table}[H]
    \centering
    \caption{Unsupervised Machine Learning Model Performance Metrics}
    \label{tab:model_performance_unsupervised}
    \scriptsize
    \setlength{\tabcolsep}{3pt}

    \begin{tabular}{llccccccc}
        \toprule
        \textbf{Model} & \textbf{Test Dataset} &
        \textbf{Accuracy} & \textbf{Precision} & \textbf{Recall} &
        \textbf{F1} & \textbf{ROC-AUC} \\
        \midrule

        \multirow[c]{2}{*}{Isolation Forest}
        & CSE-CIC-IDS2018 & 0.7409 & 0.3138 & 0.4469 & 0.3687 & 0.6947 \\ 
        & UNSW-NB15       & 0.3019 & 0.1344 & 0.8306 & 0.2314 & 0.3528  \\

        \midrule

        \multirow[c]{2}{*}{One-Class SVM}
        & CSE-CIC-IDS2018 & 0.7851 & 0.1585 & 0.0625 & 0.0897 & 0.7241 \\
        & UNSW-NB15       & 0.5948 & 0.0194 & 0.0444 & 0.0270 & 0.4193 \\

        \midrule

        \multirow[c]{2}{*}{Improved Auto-Encoder}
        & CSE-CIC-IDS2018 & 0.7129 & 0.2237 & 0.2818 & 0.2494 & 0.6534 \\
        & UNSW-NB15       & 0.5063 & 0.1844 & 0.8481 & 0.3029 & 0.7573 \\ 

        \bottomrule
    \end{tabular}
\end{table}
\FloatBarrier

Unlike the performance of supervised machine learning models, all the three unsupervised 
machine learning models generated accuracy lesser than 80\% for CSE-CIC-IDS2018 dataset 
and below 60\% while tested on UNSW-NB15 dataset. Isolation Forest showed superior performance 
while generalizing for CSE-CIC-IDS2018 dataset with F1\_Score of 0.3687 and ROC-AUC score of 
0.6947. Meanwhile, Improved Auto-Encoder outperformed other three unsupervised models while evaluating 
the UNSW-NB15 dataset with F1\_Score of 0.3029 and ROC-AUC score of 0.7573. Unlike XGBoost, which 
turned out to be the most robust model for cross-dataset generalization (considering both test datasets),
no single unsupervised machine learning models showed superior capability for cross-dataset generalization 
for network intrusion detection. One-Class SVM turned out to be the least performing unsupervised model 
for cross dataset generalization with F1\_Score of 0.0897 and 0.0270 for CSE-CIC-IDS2018 and UNSW-NB15 
datasets respectively. 
Furthermore, the results of both the supervised and unsupervised machine learning models imply that no 
standalone machine learning models have extra capability to generalize cross-dataset intrusion detection 
while evaluating on benchmark datasets considering huge distribution shifts and class imbalance in the datasets.

\section{Hypothesis Analysis}
In adddition to the results obtained after training the models to find which model shows greater capabilities in cross-dataset generalizatio 
for NIDS, we also conducted a 10 fold significance test (p-value test) for both supervised and unsupervised machine learning models to check whether 
our hypotheses were significant or not. The results showed that 5 of our hypotheses were significant. H5, H7 and H8 although cannot be determined by 
a significance test, based on the descriptive analysis (mean cross F1\_Score) we have considered H5 to be significant. However, H7 and H8 requires a 
predeclared feature intervention and mediation design; thus remains a further analysis.

\begin{table}[!ht]
    \centering
    \caption{Hypothesis Results}
    \label{tab:hypothesis-results}
    {\fontsize{9}{11}\selectfont
    \renewcommand{\arraystretch}{1.2}
    \setlength{\tabcolsep}{4pt}
    \begin{tabularx}{\textwidth}{
        >{\raggedright\arraybackslash}p{0.8cm}
        >{\raggedright\arraybackslash}p{8.5cm}
        >{\raggedright\arraybackslash}p{2.2cm}
        >{\raggedright\arraybackslash}X
    }
        \toprule
        \textbf{S.No} & \textbf{Hypothesis} & \textbf{Significance Test (p-value)} & \textbf{Remarks} \\
        \midrule

        H1 & The evaluation metrics Precision, Recall, F1\_Score and ROC-AUC for intrusion detection depend significantly on the type of model used.
        & min Holm $p = 0.35156$
        & Significant \\

        H2 & The evaluation metrics of the models for intrusion detection depend significantly on the assessment setting used for evaluation.
        & min Holm $p = 0.35156$
        & Significant \\

        H3 & Attack type signifies the performance of the models in terms of evaluation metrics.
        & min Holm $p = 0.35156$
        & Significant \\

        H4 & Cross-generalization for network intrusion detection effectiveness relies on the interaction between both the selection of models and evaluation settings.
        & min Holm $p = 0.35156$
        & Significant \\

        H5 & Robustness is not determined solely by the highest score on one transfer path.
        & Not a $p$-value hypothesis
        & Results based on mean, worst path, and variability across all cross-dataset paths \\

        H6 & Evaluation setting and attack type have significance in generalization for intrusion detection.
        & min Holm $p = 0.35156$
        & Significant \\

        H7 & Dataset features have an intermediary role between the individual variables and dependent variables for the effectiveness of intrusion detection.
        & Not estimable
        & Requires a predeclared feature intervention and mediation design. \\

        H8 & Feature selection/type mediates the link between the model and the performance of intrusion detection.
        & Not estimable
        & Requires a predeclared feature intervention and mediation design. \\

        \bottomrule
    \end{tabularx}
    }
\end{table}
\FloatBarrier
\clearpage

\begin{table}[!ht]
    \centering
    \caption{Cross-path robustness across models (Descriptive Dtatistics for H5).}
    \label{tab:cross-path-robustness}
    {\fontsize{9}{11}\selectfont
    \renewcommand{\arraystretch}{1.2}
    \setlength{\tabcolsep}{6pt}
    \begin{tabularx}{\textwidth}{
        >{\raggedright\arraybackslash}X
        >{\centering\arraybackslash}p{2.8cm}
        >{\centering\arraybackslash}p{2.8cm}
        >{\centering\arraybackslash}p{2.8cm}
        >{\centering\arraybackslash}p{2.8cm}
    }
        \toprule
        \textbf{Model} & \textbf{mean\_cross\_auc} & \textbf{worst\_cross\_auc} & \textbf{sd\_cross\_auc} & \textbf{mean\_cross\_f1} \\
        \midrule

        \multicolumn{5}{l}{\textbf{Supervised Models}} \\
        \midrule
        XGBoost                 & 0.5961 & 0.1872 & 0.1982 & 0.0883 \\
        Deep Neural Network     & 0.5806 & 0.1367 & 0.1986 & 0.1517 \\
        Random Forest           & 0.5755 & 0.1401 & 0.2284 & 0.0561 \\

        \midrule
        \multicolumn{5}{l}{\textbf{Unsupervised Models}} \\
        \midrule
        Improved Autoencoder    & 0.6470 & 0.2644 & 0.2012 & 0.2537 \\
        One-Class SVM           & 0.6074 & 0.2708 & 0.2173 & 0.1780 \\
        Isolation Forest        & 0.5975 & 0.2503 & 0.2251 & 0.2294 \\

        \bottomrule
    \end{tabularx}
    }
\end{table}
\FloatBarrier







